%% The countable atomless Boolean algebra and its Stone space.

\subsection{The countable atomless Boolean algebra and Cantor space}
\label{sec:p10-countable-atomless}

The finite-splitting and countable constructions below are proved in
\(\mathsf{ZF}\). The final correspondence for arbitrary Boolean
algebras states its additional BPIT assumption explicitly.

\begin{hbdefinition}{def:p10-atomless-countability}{Atomlessness and countability}
In a Boolean algebra \(B\), an \emph{atom} is an element \(p\ne0\)
such that \(0<b\leq p\) implies \(b=p\).
The algebra is \emph{atomless} if
\[
 \forall b\ne0\ \exists c\quad 0<c<b.
\]
It is \emph{nontrivial} when \(0\ne1\).
Here \emph{countable} means finite or countably infinite.
Every finite nontrivial Boolean algebra has an atom, so a nontrivial
countable atomless algebra is necessarily countably infinite.
The nontriviality hypothesis excludes the one-element algebra, which
is atomless vacuously and has empty Stone space.

All Boolean subalgebras contain the same \(0,1\) as their ambient
algebra.  Atoms of a finite subalgebra are minimal nonzero elements
of that subalgebra; they can split further in an atomless ambient
algebra.
\index{atomless Boolean algebra}\index{countable atomless Boolean algebra}
\end{hbdefinition}

\begin{hbtheorem}{thm:p10-atomless-extension}{proved=annex:p10.atomless.1}{Extending an isomorphism of finite Boolean subalgebras}
Let \(B,C\) be nontrivial Boolean algebras, with \(C\) atomless.
Suppose \(F\leq F'\leq B\) and \(G\leq C\) are finite subalgebras,
and \(h:F\to G\) is an isomorphism.
Then there are a finite subalgebra \(G'\) of \(C\) containing \(G\)
and an isomorphism \(h':F'\to G'\) extending \(h\).

More explicitly, if an atom \(p\) of \(F\) is the join of \(k_p\)
atoms of \(F'\), split \(h(p)\) into \(k_p\) pairwise disjoint
nonzero elements of \(C\), and map those atoms bijectively to the
chosen pieces.
\end{hbtheorem}

\begin{hbtheorem}{thm:p10-countable-atomless-uniqueness}{proved=annex:p10.atomless.2}{Uniqueness of the countable atomless Boolean algebra}
Any two nontrivial countable atomless Boolean algebras are isomorphic.
More generally, every isomorphism between finite subalgebras of two
such algebras extends to an isomorphism of the whole algebras.
\end{hbtheorem}

The extension theorem is used alternately in the two algebras:
one step includes the next element in the domain, and the next
includes the next element in the range.
This is the back-and-forth argument suggested in
\cite[Chapter~IV, \S1, Exercise~5, p.~121]{BurrisSankappanavar2012};
the annex supplies the splitting and union proofs.

\begin{hbdefinition}{def:p10-cantor-cylinder-algebra}{Cantor space and its cylinder algebra}
Put \(K=2^{\mathbb N}=\{0,1\}^{\mathbb N}\), with the product
topology of the discrete two-element space.
Coordinates are numbered \(0,1,2,\ldots\).
For \(s\in\{0,1\}^n\), its \emph{cylinder} is
\[
 [s]=\{x\in K:x|_{\{0,\ldots,n-1\}}=s\}.
\]
The cylinder of the empty word is \(K\).
Write \(x|_n\) for the initial segment of \(x\) of length \(n\).
These cylinders form a clopen basis; this product space is called
\emph{Cantor space}.
Define
\[
 \mathcal B_n
   =\left\{\ \bigcup_{s\in E}[s]:E\subseteq\{0,1\}^n\right\},
 \qquad
 \mathcal B_{\rm cyl}=\bigcup_{n\in\mathbb N}\mathcal B_n.
\]
The Boolean operations are union, intersection, and complement
relative to \(K\), with constants \(\varnothing,K\).
\index{Cantor space}\index{cylinder set}\index{Boolean algebra!cylinder algebra}
\end{hbdefinition}

\begin{hbtheorem}{thm:p10-cantor-cylinder-algebra}{proved=annex:p10.atomless.3}{Existence through clopen cylinders}
Cantor space is nonempty, compact, and Hausdorff.
Every clopen subset of \(K\) is a finite union of cylinders of
one common length, so
\[
 \operatorname{Clop}(K)=\mathcal B_{\rm cyl}.
\]
This is a nontrivial countably infinite atomless Boolean algebra.
The finite subalgebras \(\mathcal B_n\) have
\(2^{2^n}\) elements; their inclusions split every cylinder atom
\([s]\) into \([s0]\) and \([s1]\).
Their union is also the direct limit of this inclusion sequence.
Consequently there exists exactly one nontrivial countable atomless
Boolean algebra up to isomorphism.
\end{hbtheorem}

\begin{hbtheorem}{thm:p10-countable-atomless-stone}{proved=annex:p10.atomless.4}{The Stone space is Cantor space}
Let \(B\) be a nontrivial countable atomless Boolean algebra.
Choose an isomorphism
\(\alpha:B\to\mathcal B_{\rm cyl}=\operatorname{Clop}(K)\).
Then
\[
 H:K\longrightarrow X_B,\qquad
 H(x)=\{b\in B:x\in\alpha(b)\}
\]
is a homeomorphism, with
\[
 H^{-1}(\widehat b)=\alpha(b).
\]
Thus every ultrafilter of \(B\) corresponds to exactly one infinite
binary sequence.
\end{hbtheorem}

The Cantor-space connection is stated for a countably infinite free Boolean
algebra in \cite[Chapter~IV, \S4, Exercise~3, p.~140]{BurrisSankappanavar2012}.
Here existence, algebraic uniqueness, and the displayed point map
give the conclusion directly for every nontrivial countable
atomless Boolean algebra.

\begin{hbtheorem}{thm:p10-atoms-isolated}{proved=annex:p10.atomless.5}{Atoms and isolated Stone points}
Assume the Boolean prime ideal theorem, as in
Theorem~\ref{thm:p10-stone-representation}.
For every Boolean algebra \(B\), its atoms correspond bijectively
to the isolated points of \(X_B\).
The atom \(p\) corresponds to the principal ultrafilter
\[
 U_p=\{b\in B:p\leq b\},\qquad \widehat p=\{U_p\}.
\]
In particular, \(B\) is atomless exactly when \(X_B\) has no
isolated points.  If \(B\) is finite and nontrivial with \(k\)
atoms, then \(X_B\) is a discrete \(k\)-point space and
\(|B|=2^k\).
\end{hbtheorem}

This is the correspondence in
\cite[Chapter~IV, \S4, Exercise~11, p.~141]{BurrisSankappanavar2012}.
It also explains why the finite nontrivial cases cannot have
Cantor space as their Stone space: every cylinder in \(K\)
contains two smaller nonempty cylinders, so \(K\) has no isolated
points.
